\section{Finite integration with exact boundary products}
\label{virt:section}

The local charts of the preceding section must be assembled so that Stokes'
theorem gives the \emph{same} integrations on every boundary factor.  We do
this at a fixed finite cut.  The coefficient variables of the obstruction
equations serve as continuous perturbation parameters; integrating their
normalized densities removes them.  Intrinsic incidence points on the
disks identify those variables on overlaps.  The point requiring care is
that a chart may carry more variables than its neighbor.  We arrange that
every such additional variable has identically zero insertion on the
overlap, so adding it is an ordinary product operation.

Throughout this section fix one of the finite boundary-closed systems of
Proposition~\ref{local:finite}, including its relative classes, and fix an
arbitrarily large finite differentiability order.  A \emph{map} includes
its domain, boundary lifts, and original markings, but excludes auxiliary
stabilizing markings and obstruction coefficients.  All map isomorphisms
preserve these original data.  In the $p$-sector we retain the actual
cyclically ordered input slots: cyclic reindexing is a symmetry between
these labeled spaces, not an additional quotient.

\begin{theorem}[Finite integration package]\label{virt:finite-integration}
For this finite system there are oriented integration correspondences for
the stable pure $m$, directed, mixed, and interior-output $p$ sectors,
including $p_0$, with the following properties.
\begin{enumerate}[label=\textup{(\roman*)}]
\item Pullbacks of evaluation forms can be integrated in the virtual
dimension.  The pure boundary-output correspondences give differential
forms by submersive pushforward.  The interior output of a $p$-operation
may be tested against any smooth ambient form; no submersivity of that
evaluation is required.
\item Stokes' theorem holds exactly with the normalized boundary.  On a
face with distinguished nodes its integration is the product of the
component prescriptions, using an evaluation fiber product over $L$ at
a smooth node and a product over the fixed intersection generator at a
jump node.  These identifications agree on iterated corners.
\item The prescriptions are equivariant under cyclic reindexing of $p$
inputs, with the determinant and form-permutation signs.  Only actual
finite map isotropy contributes an orbifold integration factor.
\item A regular zero-dimensional type with empty boundary and trivial
isotropy has its ordinary signed count.  In particular this applies to
the rigid directed spaces and to the low-area rigid comparison triangles
of the preceding section.  The constant $p_{1,0}$ correspondence is
ordinary integration on $L$, with a fixed orientation sign.
\end{enumerate}
All assertions hold coefficientwise after multiplying each relative-class
term by its energy monomial and, when defined, by
$\exp(tB_{\rm rel})$.  The choices need only be made at the specified
finite cut.  Operation signs and the unstable de Rham differential are
specified in the next section.
\end{theorem}

We prove the theorem by constructing and comparing the integrations.
The framework for disk obstruction charts in
\cite[Theorem~7.1 and Condition~7.11]{fooo-i} and the componentwise
framework of \cite[Theorems~2.5, 5.3--5.4 and Remark~6.1]{fooo-ii}
motivate the construction.  The intrinsic probes, the $p_0$ treatment,
and the precise overlap integrations below are part of our argument.

\subsection{Intrinsic incidence probes}

Write $\mathbb D$ for the open unit disk, equipped with its hyperbolic
metric.  A \emph{probe} consists of the following finite data:
\begin{enumerate}[label=\textup{(\alph*)}]
\item two ordered codimension-two target slices $S_1,S_2$, with compact
subpatches $K_i\subset\operatorname{int}(S_i)$;
\item a positive lower threshold for the least singular value of the
normal incidence derivative at each point, computed using the hyperbolic
metric on the disk, and bounds
$0<d_-\le d_+<\infty$ for the distance between the two points;
\item an oriented, even-dimensional real vector space $E$ and a linear
map from $E$ to smooth $TT$-valued real one-forms compactly supported in
$\mathbb D$, using the fixed parallelization of $TT$.  The map has at
least one dummy coordinate with zero image.
\end{enumerate}
Choose a compact set $K\Subset\mathbb D$ containing all these supports,
the origin, and every possible normalized second point at distance in
$[d_-,d_+]$ from the origin.  Choose
$0<\epsilon<\operatorname{dist}(K,\partial\mathbb D)$.

An occurrence on an actual map $\mathbf u$ is an ordered pair of distinct
interior points $(w_1,w_2)$ of one normalized disk component such that
$u(w_i)\in K_i$ and all the closed thresholds hold.  There is a unique
disk biholomorphism
\[
 \varphi:\mathbb D\longrightarrow D_v,
 \qquad \varphi(0)=w_1,\qquad
 \varphi(r)=w_2\quad\text{for some }r\in(0,1).
\]
If the original interior output mark belongs to this component, require
its $\varphi$-coordinate to be at distance at least $\epsilon$ from $K$.
There is no restriction involving the order or position of the boundary
special points.  We use every occurrence on the representative map,
without quotienting the occurrence set by automorphisms.

The probe also carries a smooth scalar cutoff $0\le\chi\le1$ along
transverse incidence branches.  Its support lies within strictly stronger
versions of the closed conditions just imposed.  For instance, choose
bumps on the full slices supported inside $K_i$, a distance bump supported
inside $(d_-,d_+)$, and smooth functions of the normal derivative matrices
supported where the least singular values exceed the selected
thresholds.  Such derivative cutoffs exist on the finite-dimensional
space of matrices; they can be made invariant under orthogonal changes
of frames.  For the interior mark use a smooth function on the normalized
disk, zero on a neighborhood of $K$ and equal to one near the circle,
with support inside the avoidance condition.  This factor is defined to
be one if the component has no original interior mark.  Given a desired
transverse pair with strict margins, these choices can be made with
$\chi=1$ near that pair.

\begin{lemma}[Finiteness and continuation of probes]\label{virt:probes}
For a finite list of probes, every actual stable map has finitely many
occurrences.  If actual maps converge stably, each sequence of occurrences
has a subsequence converging, in a whole-disk normalization, to an
occurrence on one disk of the limiting map.  Each limiting occurrence
continues uniquely to nearby maps, including sufficiently close ambient
maps in stabilized plumbing neighborhoods.  The continuations are
equivariant under map isomorphisms.
\end{lemma}

\begin{proof}
There is a positive lower bound for the area in a hyperbolic radius-one
ball about the first incidence point.  Otherwise recenter a sequence of
such balls by disk isometries.  The harmonic-measure coefficients have
uniform bounds of every required order on compact interior balls; only
finitely many equation data occur at this cut.  Interior small-energy
regularity would force the derivative at the center to tend to zero,
contrary to the positive lower incidence threshold.  This argument is
applied here to actual, unforced maps.

In a stably convergent sequence, the energy identity allows compact
interior cores of the nonconstant limiting disks to carry all but less
than this lower bound of the sequence's energy.  A radius-one ball at an
occurrence must therefore meet one of these cores.  The component
reparametrizations are restrictions of whole-disk biholomorphisms, as in
the boundary plumbing construction.  In the normalization of the core
that it meets, the first point belongs to a fixed compact interior set;
the upper bound $d_+$ puts the second in another such set.  After passing
to a subsequence both converge there.  Smooth core convergence preserves
the closed slice, derivative, and distance conditions.  If the limiting
component carries the original interior mark, its coordinate converges
too, preserving mark avoidance.  We obtain an occurrence on a single
limiting disk.

The derivative of the incidence equation with respect to the position
$w_i$ is an isomorphism onto the two-dimensional normal space of $S_i$.
The implicit function theorem, using the full slice rather than its
compact subpatch, therefore continues the point uniquely under
sufficiently small $C^2$ core perturbations and domain deformations.
The same argument applies simultaneously to the ordered pair.  For a
fixed map, an infinite occurrence set would have an accumulation point
by the preceding compactness argument, contradicting this local
uniqueness.  Thus the set is finite.  A sequence of occurrences on nearby
actual maps that was not among the continued branches would likewise
contradict compactness and uniqueness.  Distinct continued branches
remain distinct after shrinking.

Finally, all conditions and the normalization are invariant under disk
biholomorphisms.  Changing the core identifications continues the same
incidence branches; central finite automorphisms merely permute them.
This proves the equivariance assertion.
\end{proof}

For a chart centered at $\mathbf u$, let $O_{\mathbf u}$ be the occurrences
at its center, and continue this fixed finite set over the chart
neighborhood.  \emph{Do not retest the closed conditions} at a nearby
map $\mathbf x$.  Set
\[
 E_{\mathbf u}=\bigoplus_{o\in O_{\mathbf u}}E_o.
\]
Transport the one-forms for $E_o$ by its actual normalization
$\varphi_o(\mathbf x)$ and project to $(0,1)$-type for the equation on
$\mathbf x$.  Denote the resulting insertion by $I_o(\mathbf x)$.
The coefficient space is a formal direct sum: the insertion maps need
not be injective or have independent images.

These insertions satisfy the smooth core-transport hypotheses of
Proposition~\ref{local:charts}.  Incidence positions are recovered smoothly on
the cores.  Whole-disk normalizations commute with the projective
insertions in real boundary plumbing, and their expressions on retained
cores have no singular scaling.  Projection to the varying $(0,1)$-type
has the length-derivative estimates established there.  The derivative
and mark cutoffs are smooth finite-order core observables.  If the
original interior mark enters the disk from a different component under
smoothing, its coordinate tends to the unit circle, so its avoidance
factor is identically one near the broken configuration.  Thus on every
normalized face both the insertions and their cutoff factors restrict
\emph{exactly} to the component rules.

\subsection{A finite family of regular thickenings}

At each actual map choose two distinct immersive interior points on each
nonconstant disk, transverse slices through their images, and a probe
with strict margins at this pair.  Its compactly supported test forms
can be chosen to span the finite-dimensional cokernel of the fixed-domain
operator.  On a pure $m$ component choose them to span the cokernel of
the operator augmented by its boundary output evaluation.  Compact
interior tests away from the original interior mark suffice by
Proposition~\ref{local:charts}.  Enlarge the coefficient space to even dimension
and include a dummy zero direction.  Include a probe also on a regular
nonconstant component, since it will be useful for isotropy.

Constant components require no such tests.  A pure constant component
has the regular real-disk operator with freely varying value in $L$.
A constant mixed component has only an $X,L$ pair, two opposite jumps
at one intersection, and possibly removable smooth $L$-marks.  Its
map operator is the invertible transverse strip operator.  A stable
constant directed component cannot occur: three distinct boundary labels
have no common point, whereas a two-ended constant strip is unstable.

The chosen active branches persist on nearby actual maps.  In particular,
when several components join, the inverse construction of
Proposition~\ref{local:charts} uses their continuing probe blocks.  All smooth
matching equations are solved using pure child output submersivity;
one may still prescribe the pure parent output.  This proves the
required fixed-domain augmentations for each joined component.  A nearby
nonconstant joined component contains a continuing probe: a cluster of
constant components cannot become a nonconstant actual component with
energy tending to zero, by the area gap.  Compactness now gives a finite
cover of all actual moduli at the cut by such neighborhoods.  Take the
union of their probe lists and allow every probe on every disk.
Additional coefficient directions preserve the augmentations.

For small coefficients we use the space of pairs satisfying
\begin{equation}\label{virt:thickening}
 \bar\partial_J\mathbf x
   =\sum_{o\in O_{\mathbf u}}
        \chi_o(\mathbf x)I_o(\mathbf x)e_o,
 \qquad e_o\in E_o.
\end{equation}
After auxiliary stabilization and quotient by actual map isomorphisms,
Proposition~\ref{local:charts} gives a regular orbifold with corners near
$(\mathbf u,0)$.  Its dimension is the virtual dimension plus
$\dim E_{\mathbf u}$.  Pure boundary output is submersive on every
stratum.  Auxiliary incidence marks can be chosen in automorphism orbits
away from the compact obstruction supports: nonconstant maps have
immersive points in any remaining interior collar, by unique
continuation.  The auxiliary mark variables and slice constraints cancel
in the tangent determinant.  Orient the thickening by the virtual
Fredholm and marked-domain determinant, followed by the oriented
coefficient directions.

Here is the isotropy check, including the sector without boundary inputs.
For $p$, root the disk tree at its interior mark.  If a nonconstant root
has a nonidentity finite stabilizer, an ordered pair of distinct interior
anchors cannot be fixed: a disk automorphism fixing both is the identity.
The stabilizer consequently permutes probe occurrences nontrivially.
Their dummy coefficients, which are tangent kernel directions at zero,
detect this permutation even after quotient by infinitesimal
reparametrizations.  A constant root has at least one boundary special
point by stability; with only one, its stabilizer is trivial.  A
nontrivial rotation about the interior mark must permute at least two
boundary special points nontrivially and acts effectively on their
angular configurations modulo simultaneous rotation.  Indeed the tangent
space of these configurations is $\R^r/\R(1,\ldots,1)$ for $r\ge2$;
no nonidentity permutation induced by a finite rotation acts trivially
on it.  Fixed-domain augmented surjectivity allows all these domain
variations in the thickening.  An automorphism identical on the root is
identical down the tree, because a finite disk automorphism fixing a
boundary parent point is the identity.  The same last argument makes
boundary-rooted $m$ stabilizers trivial.  Thus the charts can be taken
effective.  For an unbroken nonconstant $p_0$ map the two interior anchors
provide exactly the stabilization and effective action just described.

\subsection{Comparison of charts without changing the equation}

\begin{lemma}[Joint charts and inactive coefficients]\label{virt:joint-charts}
The thickening neighborhoods can be shrunk with the following properties.
Every incidence branch with potentially nonzero cutoff is tracked there.
On an overlap, branches common to two charts have identical coefficient
coordinates and insertions.  Every branch belonging to only one list has
identically zero cutoff near the overlap.  After adding these inactive
blocks to the other list, the two equations define the same local space
of pairs, with the same evaluations, isotropy, and determinant orientation.
The comparison is smooth to the prescribed finite order, including on
faces with any specified nodes retained.
\end{lemma}

\begin{proof}
First exclude untracked active branches for ambient maps, not merely
actual ones.  Choose fixed compact cores and sufficiently small gluing
neighborhoods.  Outside the cores the derivative on hyperbolic unit
vectors can be required to be smaller than every probe's lower
threshold.  At ordinary boundary arcs this follows from smooth boundary
control and the blowup of the hyperbolic metric; at long necks and jump
ends it follows from the decay estimates in Lemma~\ref{local:decay},
where hyperbolic unit vectors have bounded strip-coordinate norm.
These restrictions are valid also for the small-coefficient solutions,
whose obstruction terms remain on the specified cores.  Consequently
any pair in the inner cutoff supports must lie on the continued compact
cores.  Core closeness and the implicit uniqueness in
Lemma~\ref{virt:probes} identify it with a tracked branch.  The strict
margin between the cutoff support and the closed occurrence conditions
is essential: a branch failing the latter at the center stays outside
the former on a neighborhood.

We next make precise that the underlying-map comparisons used here have
no extra multiplicity.  The core convergence in these neighborhoods,
together with arbitrarily small oscillation in the omitted necks and
ends, has unique stable-map limits up to isomorphism.  To prove this,
suppose a sequence admits two such limits.  Stabilize using the markings
for both limits and pass to a common stable marked-domain limit.  Extra
comparison marks may be displaced slightly to avoid coincidences.
Forgetting one or the other marking system recovers the two given domain
limits.  A component contracted in one description has constant limiting
map, by core convergence or the small oscillation on its omitted region.
The maps on every retained component agree in the common refinement.
After forgetting all auxiliary marks and contracting unstable constant
components, both limits therefore give the same stable map.  This uses
no holomorphic equation on the retained cores.

It follows in particular that, after shrinking about a center, lifts of
the same map to its slice gauge differ only by its central finite
automorphisms.  Otherwise compare a sequence of conflicting gauge lifts
by the preceding marked convergence.  Their limiting identifications
are central automorphisms, which preserve the chosen slice systems with
their induced relabelings.  Transverse incidence then makes the continued
marks exactly identical up to that relabeling, a contradiction.  This
also shows that underlying isomorphisms lift to the gauges.  Stability is
preserved nearby, either by the original domain specials or by the
persistent immersive cores of nonconstant components.

On an overlap take the union of the two continued branch lists, identifying
coinciding branches by transverse uniqueness.  A branch missing from
one list is outside its inner cutoff support throughout a smaller
comparison neighborhood by the first paragraph; hence its contribution
to Equation~\eqref{virt:thickening} is identically zero there.  Adding its
coefficient block thus takes a product with $E_o$, rather than changing
the equation.  Redundancy among the active insertions does not affect
this assertion: coefficients remain part of the pair and are never
quotiented by the kernel of their insertion.

For smoothness recover both sets of stabilizing marks by their transverse
incidences, or compare through a third marking system disjoint from them
at the point in question.  Each stabilization retains every component
of that point.  Adding and forgetting comparison marks therefore
contracts no component there.  Every retained node represents the same
single domain divisor on the two sides, so its real plumbing parameters
differ by a smooth positive unit.  The inverse-length changes are smooth
by Proposition~\ref{local:charts}; there is no multiplication of plumbing
parameters along a contracted chain.  This argument also applies on a
normalized face, keeping its distinguished nodes.  Core observables,
slice positions, and kernel coordinates compare smoothly, and the
weighted neighborhood conditions are preserved after shrinking by the
same decay estimates.  It applies with small nonzero coefficients,
since the padded equations are literally identical.

These comparisons retain the map isomorphisms, acting on the branch
variables as well.  The auxiliary mark and constraint determinants
cancel, and adjoining a block appends its oriented determinant.
Permuting entire blocks preserves orientation because their ranks are
even.  This proves all the claims.
\end{proof}

\subsection{Buffered partitions depending only on maps}

Choose for each probe a smooth compactly supported top-form density
$\Xi^\hbar$ on its oriented coefficient space, of integral one and
supported in $|e|<\hbar$.  The same density is used at every occurrence
of that probe.  In a chart let
\begin{equation}\label{virt:density}
 \Xi_{\mathbf u}^\hbar
       =\bigwedge_{o\in O_{\mathbf u}}\Xi_o^\hbar(e_o).
\end{equation}
It is closed, and its definition is unaffected by permutation of the
even-dimensional blocks.  Inactive blocks still carry their full
densities.

We must localize on maps without selecting only some coefficient
representatives of a map.  The following elementary consequence of the
chart construction supplies this point.

\begin{lemma}[Buffered partitions]\label{virt:buffered-partitions}
There are finitely many trimmed thickening charts $V_\alpha$, nonnegative
map functions $f_\alpha$, and a common sufficiently small $\hbar>0$,
such that near all density supports:
\begin{enumerate}[label=\textup{(\roman*)}]
\item the zero extensions of the $f_\alpha$ pull back smoothly to every
chart, and $\sum_\alpha f_\alpha>0$;
\item each $f_\alpha$ has support, with a buffer, strictly inside its
chart, relative to the actual nodal corners;
\item if a map is in that support, the chart represents every assignment
at the common small support scale that solves its equation for that map.
This includes every free choice of coefficients in any blocks made
inactive during an overlap comparison.
\end{enumerate}
Thus $\lambda_\alpha=f_\alpha/\sum_\beta f_\beta$ is a partition on
the map projections of the relevant pairs.  It is independent of which
small coefficient assignment represents the map.
\end{lemma}

\begin{proof}
Around each center choose nested coordinate neighborhoods with compact
closures inside a larger thickening chart, and trim inside the larger
one.  Here compactness refers to the finite-dimensional chart, with its
actual corner faces retained.  Its local coordinates can be chosen from
the coefficient coordinates and finitely many smooth functionals of the
map on retained cores, together with stabilized domain variables and
inverse neck lengths.  Indeed these functionals separate the augmented
kernel: after fixing the coefficients, a homogeneous kernel section
vanishing on open cores vanishes by elliptic uniqueness.  Finitely many
tests suffice on this finite-dimensional kernel, and the inverse
function theorem supplies local coordinates.

Consequently bounds on finitely many map and domain observables keep a
solution in a smaller coordinate neighborhood uniformly for all
sufficiently small coefficients.  Take a smooth bump of these
observables, positive at the center and supported with a strict interior
buffer.  Average over the central finite group.  The observables use
continued slice gauges, whose only ambiguity is that group by
Lemma~\ref{virt:joint-charts}; the averaged function is therefore a
function of the map.  These bumps do not depend on its coefficients.
Select finitely many so that their positive neighborhoods cover the
actual compact locus, and include their cyclic translates in the
$p$-sector.

We justify the required uniformity before fixing the support radius.
Consider a sequence of pairs in the trimmed closures whose coefficients
tend to zero.  There is a subsequence converging in a larger chart to
$(\mathbf v,0)$, where $\mathbf v$ is actual.  Every applicable open
gauge neighborhood of $\mathbf v$ then contains the maps in that
sequence.  Transfer their common branch coefficients to such a gauge
and choose arbitrary coefficients tending to zero in its inactive
blocks.  The local chart parametrizes all small solutions of its
equation, so every resulting pair tends to $(\mathbf v,0)$ in that chart.
Uniqueness of the map limit and slice positions excludes any other
nearby lift.  Some $f_\alpha$ is positive at $\mathbf v$; thus failure
of positivity at every sufficiently small coefficient scale would
contradict this subsequence argument.

The same argument proves the all-assignments assertion.  If it failed
at arbitrarily small scales, take an offending map in a projected cutoff
support and an alternative assignment at that scale solving the same
equation.  A lift of the map in the originating cutoff support lies in
its compact buffered region.  Pass to its actual limit.  The alternative
assignment tends to zero; after branch transfer and arbitrary padding
it must lie near $(\mathbf v,0)$ in the originating chart, by its local
completeness.  The strict buffer contradicts its asserted exclusion.
This reasoning includes the inactive variables and any kernel among
the active insertions.  Finiteness of the chart collection permits one
uniform scale.  One may first obtain these conclusions for coefficients
bounded by $2\hbar$, and then put the densities inside the $\hbar$ balls.

Finally consider smoothness of the zero extensions at this fixed small
scale.  Near a projected cutoff support, the compact buffer and
uniqueness of underlying-map limits put every comparison in the same
open slice gauge.  The map observables there compare smoothly by
Lemma~\ref{virt:joint-charts}, and the all-assignments property permits
every small coefficient transfer needed for that comparison.  Away
from the projected support the cutoff is identically zero on a
neighborhood of the density support.  Its extension is therefore
smooth in either case.  The same buffer applies to its derivatives.
Normalizing the sum now gives the asserted partition.
\end{proof}

This construction uses no global parameter space of constant rank.
Only the local pairs and the map functions in
Lemma~\ref{virt:buffered-partitions} enter the integration.

\subsection{Patching, Stokes, and boundary composition}

For a product $\alpha$ of pulled-back evaluation forms, in its specified
order, define
\begin{equation}\label{virt:integral}
 \mathcal I_{\mathcal M}(\alpha)
   =\sum_\gamma\int_{V_\gamma}
           \lambda_\gamma\,\alpha\wedge\Xi_\gamma^\hbar.
\end{equation}
Integrals are ordinary orbifold integrals: in a local finite quotient
they are divided by its actual group order.  The density is put last.
The integral vanishes unless the remaining form degree is the virtual
dimension.  Compact support in the coefficient directions and the
buffered map support make every integral compactly supported away from
the artificial trimming boundary.  At a pure boundary output the same
formula, using submersive fiber integration, gives a differential form.
We characterize its order convention by placing the output test form
first.  Submersivity on every stratum ensures that its boundary
pushforwards are defined by the same rule.

\begin{lemma}[Exact patching and Stokes]\label{virt:patching}
Equation~\eqref{virt:integral} is independent of the buffered partition
and of the slice gauges used to compute it, for the chosen probes and
densities.  It is unchanged by adjoining inactive blocks.  It satisfies
\[
 \mathcal I_{\mathcal M}(d\alpha)
       =\mathcal I_{\partial^{\rm norm}\mathcal M}(\alpha),
\]
where the right side has the induced boundary orientation.  Its
prescription on a normalized face is exactly the product prescription
from its component moduli, with the corresponding matching conditions.
\end{lemma}

\begin{proof}
On a buffered overlap identify common branch coordinates and take the
union $O$ of the branch lists.  By Lemma~\ref{virt:joint-charts} its pair
space is locally the first chart times the missing inactive blocks, and
also the second chart times its missing inactive blocks.  Over every
map in the cutoff support these products contain the full coefficient
supports by Lemma~\ref{virt:buffered-partitions}.  All map forms and
cutoffs pull back from the same map.  Fubini therefore gives, for each
missing block,
\[
 \int_{E_o}\Xi_o^\hbar=1,
\]
and identifies the two localized integrals with the single integral
over the joint list $O$.  This comparison retains finite automorphism
quotients: locally both descriptions are quotients of the same slice
gauge with the same branch action.  It therefore introduces no extra
incidence multiplicity.  Branch continuation may permute entire blocks;
their even dimensions preserve the density and determinant orientation.
One can localize this argument on finitely many map neighborhoods with
buffers of the preceding kind, so it proves equality on the whole
overlap.  No quotient is taken by redundancy among active coefficients.

Multiplying by a second map partition and using these overlap equalities
proves partition independence.  It also proves the cancellation needed
for Stokes.  Each $\Xi_\gamma^\hbar$ is closed.  Apply ordinary Stokes
to $\lambda_\gamma\alpha\wedge\Xi_\gamma^\hbar$ and sum.  The
terms involving $d\lambda_\gamma$ can be compared on joint charts as
just above, and there they sum to zero because
$\sum_\gamma d\lambda_\gamma=0$.  There are no trimming-boundary
terms.  The only terms left are the actual normalized faces, each with
its induced orientation.  This proves the displayed identity.  The
same argument with an output test form proves its pushforward version,
including the usual signs from differentiating that form.

Now fix a face with a distinguished breaking.  The actual maps on it
are precisely the component maps with the indicated matching.  Probe
occurrences on its normalized disks, their continued insertions, and
their cutoff factors are the component prescriptions already proved
above.  Hence its pair equation is the product of the component pair
equations, with a value match at a smooth node or the fixed generator
at a jump.  A face chart and product charts may retain different
inactive branches; adding their joint lists gives the identical
equations and identical product densities.

For completeness the neighborhoods on the two sides cover the same
portion of these equations at one sufficiently small common support
scale.  An escaping sequence of component products with scale tending
to zero has, by component compactness, an actual product limit.  This
is a covered actual nodal map.  Gauge transfer and the local completeness
of the face chart then put all those assignments in its neighborhood.
Conversely a sequence in a covered face tending to zero coefficients
breaks into covered actual component limits and enters their product
neighborhoods.  The buffered all-assignments argument proves the
assertion simultaneously for all the finitely many faces and corners.

Partition independence now allows refinement to the product of the
component partitions.  Ordinary fiber integration gives exactly their
composition.  At each smooth node it is the pure child boundary-output
submersion that makes the fiber product regular and supplies this Fubini
identity; at jump nodes the matching base is a point.  Rooting a $p$
tree at the interior mark gives the same description with $m$ branches
attached in their linear input orders.  Repeating with several
distinguished nodes gives the same common component equation, proving
compatibility on all iterated corners.
\end{proof}

\begin{proof}[Completion of Theorem~\ref{virt:finite-integration}]
The construction and Lemma~\ref{virt:patching} give the integrations,
submersive outputs, and exact boundary products.  The spin and capping
determinants orient the virtual problems; appending the even coefficient
spaces preserves the determinant gluing comparisons.  The next section
computes their signs.  Our probes ignore boundary slot positions and
are therefore intrinsic under cyclic reindexing of $p$ inputs.  Pulling
back a partition by that symmetry changes no integral by partition
independence.  This proves cyclic equivariance on the labeled spaces,
with no division by the order of the reindexing group.  The effective
isotropy discussion above includes $p_0$.

For a regular rigid actual map without isotropy, the ordinary linearized
problem is invertible after its domain slice.  For every sufficiently
small choice of the probe coefficients there is a unique nearby
solution, with the same orientation sign.  Its parameter integral is
that sign times the product of the density integrals, hence that sign.
Compactness excludes new solutions away from these neighborhoods as the
support scale shrinks.  An empty actual moduli space contributes zero
for the same reason.  This proves agreement with ordinary signed counts
for a regular zero-dimensional type with empty boundary.  Directed
rigid spaces have empty boundary because all their component problems
are regular and match over points: every broken configuration has
negative expected dimension.  The low-area no-extra-mark comparison
triangles have only the same regular strip and triangle sectors, and
the energy bound excludes pure disk bubbling.  They have the same
property.  The ordered jumps exclude disk automorphisms; a nonconstant
strip has no stabilizing translation.

The constant disk with one interior and one boundary mark has no domain
moduli or automorphisms.  Its moduli space is $L$, its two evaluations
are $i_L$ and the identity, and it has no probe occurrences.  Thus the
raw $p_{1,0}$ pairing is the ordinary integral of the pulled-back ambient
test form followed by the input form on $L$, up to the fixed global
operation sign selected in the next section.  In particular its scalar
factor is nonzero and its weight is one.

Finally, energy and the relative $B$-integral are constant on the retained
relative-class refinements, also for the ambient maps in the
thickenings: they are homotopy invariants with the specified boundary
conditions and fixed jumps.  They add under every component splitting.
Their monomial and exponential factors therefore multiply in the
boundary products and have zero exterior derivative on each chart.
Stokes and all overlap identities consequently hold coefficientwise in
$t$.  The finite cut bounds the number of compositions needed later,
independently of Taylor order; the initial differentiability can be
chosen high enough for all the required pushforwards and primitives.
\end{proof}
